\documentclass[11pt,a4paper]{article}
\usepackage[margin=2.5cm]{geometry}
\usepackage{amsmath,amssymb,amsthm}
\usepackage{graphicx}
\usepackage{booktabs}
\usepackage{multirow}
\usepackage{algorithm}
\usepackage{algorithmic}
\usepackage{hyperref}
\usepackage[round]{natbib}
\usepackage{subcaption}

\newtheorem{definition}{Definition}
\newtheorem{proposition}{Proposition}
\newtheorem{finding}{Finding}
\newcommand{\eg}{e.g.,\ }
\newcommand{\ie}{i.e.,\ }

\title{When Not to Reconfigure: Sequential Expert Routing for\\
Dynamic Warehouse Slotting under Non-Stationary Demand}

\author{%
  Anonymous\\
  \texttt{(under review)}
}
\date{}

\begin{document}
\maketitle

\begin{abstract}
Dynamic slotting in fast-moving consumer goods (FMCG) warehouses is
typically optimized per-period or over short rolling horizons, treating
each re-slotting decision as independent of its predecessors. We show that
this myopic approach can be globally suboptimal when physical
reconfiguration costs create inter-temporal coupling: a locally beneficial
reconfiguration made before a demand or cost shock can lock the warehouse
into expensive future adjustments. We formalize this as the
\emph{Dynamic Warehouse Expert Routing Problem} (DWERP), a
warehouse-specific instance of switching-cost online optimization in which
a library of operations research (OR) experts proposes slotting actions
each period, and the decision is which expert to follow. Through
controlled experiments validated against exact enumeration on small
instances, we characterize the conditions under which myopic expert
selection fails: inter-temporal traps concentrate at intermediate lead
times before cost shocks (1.26--1.69\% normalized trap gain), the
full-information opportunity is modest ($\sim$1--5\%) but reproducible, and
traps are entirely absent without cost shocks. Counter-intuitively,
deployable receding-horizon policies with crude internal cost models
capture most of this opportunity not by forecasting better, but by being
conservative---model imprecision acts as implicit stickiness that avoids
over-reconfiguration. We empirically demonstrate that algorithm-switch
counts are an inadequate surrogate for physical reconfiguration cost
(hidden reconfiguration rate 17--69\%), and that value-per-move, not move
frequency, is the correct selective-reconfiguration metric. On real
warehouse data (WEPA, 3 months, 411{,}830 orders), no inter-temporal trap
manifests under natural conditions, bounding trap-aware expert routing to
genuinely non-stationary scenarios such as promotion seasons and
product-line transitions.
\end{abstract}

\section{Introduction}\label{sec:intro}

Warehouse slotting---the assignment of stock-keeping units (SKUs) to
storage locations---is a classical operations research problem whose
solution directly determines order-picking travel distance, replenishment
frequency, and congestion. In FMCG distribution centers, demand is
non-stationary (promotions, seasonal shifts, product lifecycle turnover),
so the optimal layout changes over time. \emph{Dynamic slotting} methods
address this by re-optimizing the layout periodically or over rolling
horizons.

The prevailing approach treats each re-optimization as an independent
instance: given the current demand forecast, solve for the best layout.
This is \emph{myopic} in the technical sense---it minimizes the current
period's cost without considering how today's reconfiguration changes
tomorrow's options. When physical reconfiguration is free, myopia is
harmless. When it is costly, a tension emerges: a reconfiguration that is
locally beneficial (reduces current-period picking cost) may be globally
harmful if it must be partially undone at high cost when demand shifts
again.

We call this phenomenon an \emph{inter-temporal reconfiguration trap}
(Section~\ref{sec:trap}). The core question of this paper is:

\begin{quote}
\textbf{When should a warehouse defer a locally optimal physical
reconfiguration, and how much of the resulting full-information
opportunity can a deployable policy capture?}
\end{quote}

This question sits at the intersection of three literatures: metrical task
systems and smoothed online optimization (switching costs),
warehouse online reoptimization, and algorithm selection
(Section~\ref{sec:related}). Our contribution is not a new sequential
decision theory---MTS/SOCO provides that---but a warehouse-specific
instantiation with discrete physical layout states, heterogeneous OR
experts as candidate policies, real reconfiguration costs, non-stationary
demand, and a detailed empirical characterization of when myopic
optimization fails and when it does not.

\subsection{Contributions}

\begin{enumerate}
\item \textbf{Formulation.} We define DWERP as
$\pi(S_t) \to E_t \to A_t$, $S_{t+1} = f(S_t, A_t, \xi_{t+1})$, with
transition cost $C_{\text{trans}} = d(L_t, L_{t+1})$ measured in physical
layout distance, not algorithm-switch indicators
(Section~\ref{sec:formulation}).

\item \textbf{Trap characterization.} Through controlled experiments
validated against exhaustive enumeration (beam search = exact optimum on
small instances, gap 0.00\%), we show: (i) traps exist (constructive
evidence: accepting a 1.6\% current-period loss avoids 50 costly moves,
yielding 5.34\% long-term gain); (ii) they are rare-but-costly (1/12 seeds
material, insurance-shaped distribution); (iii) they concentrate at
intermediate lead times ($\Delta t = 1$ day between demand transition and
cost shock); (iv) they vanish without cost shocks
(Sections~\ref{sec:trap},~\ref{sec:controlled}).

\item \textbf{The deployment paradox.} Deployable receding-horizon
policies with crude internal cost models capture most of the
full-information opportunity---not through better forecasting, but through
implicit conservatism. Model fidelity does not gate capture
(Section~\ref{sec:deployment}).

\item \textbf{Transition cost surrogate failure.} Algorithm-switch counts
$\mathbb{1}[E_t \neq E_{t-1}]$ are an inadequate proxy for physical
reconfiguration: hidden reconfiguration (same expert name, massively
different layout) occurs in 17--69\% of periods
(Section~\ref{sec:transition}).

\item \textbf{Real-data boundary.} On WEPA real warehouse data (411{,}830
orders, 3 months), no trap manifests under natural or moderately stressed
conditions---bounding applicability to genuinely non-stationary events
(Section~\ref{sec:real}).
\end{enumerate}

\section{Related Work}\label{sec:related}

\subsection{Metrical Task Systems and Switching-Cost Online Optimization}

Metrical task systems (MTS)~\citep{borodin1992metrical} model an online
decision-maker choosing a state each period from a metric space, paying
both service costs and movement costs. Smoothed online convex optimization
(SOCO)~\citep{goel2009smoothed} extends this with continuous action
spaces. Learning-augmented MTS uses predictions to improve competitive
ratios~\citep{bubeck2021mtses}. Our DWERP is a warehouse-specific MTS with
discrete layout states, but we do not claim new online-algorithm theory;
rather, we instantiate MTS structure in a domain with hard capacity
constraints, heterogeneous candidate policies, and physical
reconfiguration costs that violate clean metric assumptions
(Section~\ref{sec:transition}).

\subsection{Warehouse Online Reoptimization}

Recent work studies online reoptimization in picking
operations~\citep{lorenz2024picking}: simple re-solving of batching and
routing problems as orders arrive dynamically performs surprisingly well.
Related work on anticipation~\citep{lorenz2024anticipation} analyzes
strategic waiting by pickers and finds it can be counterproductive. Our
\emph{reconfiguration deferral} is fundamentally different from
\emph{picker waiting}: we defer physical warehouse state changes (pallet
moves, pick-face reassignments), not labor scheduling decisions. The two
share the inter-temporal structure but have different cost models and
decision variables.

\subsection{Dynamic Slotting and Algorithm Selection}

Dynamic storage location assignment has evolved from turnover-based
ABC~\citep{heskett1963cube} through affinity-based
methods~\citep{kofler2014affinity, mirzaei2022correlated} to DRL-based
policies~\citep{waubert2022drl}. Algorithm selection---choosing which
solver to use per instance---is studied in
OR~\citep{kotsialos2024selection} and hyper-heuristics.
Closest to our setting, contextual bandit put-away
policies~\citep{yuan2024contextual} and future-order-driven
repositioning~\citep{karimi2026keeping} adapt storage decisions to
changing demand. None of these treat the \emph{stateful reconfiguration
trajectory} as the decision object with explicit switching costs; they
optimize per-instance or per-episode independently.

\section{Problem Formulation}\label{sec:formulation}

\subsection{DWERP}

Let $t = 1, \ldots, T$ index planning periods (days or shifts). The
warehouse state is
\begin{equation}
S_t = (\mathcal{O}_t, D_t, F_t, I_t, L_t, R_t, \mathcal{C})
\end{equation}
where $\mathcal{O}_t$ = observed orders, $D_t$ = demand history, $F_t$ =
forecast, $I_t$ = inventory, $L_t$ = physical layout (SKU$\to$location
assignment), $R_t$ = resources (labor, equipment, move-cost schedule), and
$\mathcal{C}$ = hard constraints (capacity, zone, FEFO).

A library $\mathcal{E} = \{E_1, \ldots, E_K\}$ of OR experts is available.
Each expert $E_i$ is a \emph{policy}: it reads the state and outputs a new
layout $L'_t = E_i(S_t)$ together with a cost decomposition and feasibility
check.

The routing policy selects one expert per period:
\begin{equation}
\pi(S_t) \to E_t, \quad A_t = E_t(S_t), \quad S_{t+1} = f(S_t, A_t, \xi_{t+1})
\end{equation}
where $\xi_{t+1}$ is the next period's demand/cost realization. The
long-run objective is:
\begin{equation}\label{eq:objective}
\min_{\pi} \sum_{t=1}^{T} \Big[
\underbrace{C_{\text{op}}(S_t, A_t)}_{\text{picking}} +
\underbrace{C_{\text{trans}}(L_{t-1}, L_t)}_{\text{reconfiguration}}
\Big]
\end{equation}

\subsection{Transition Cost}

We define $C_{\text{trans}} = d(L_{t-1}, L_t)$ as a function of physical
layout distance---not the algorithm-switch indicator
$\mathbb{1}[E_t \neq E_{t-1}]$. In Section~\ref{sec:transition} we
empirically justify this choice. Our primary distance proxy is the number
of SKU relocations (Hamming distance over the layout map):
\begin{equation}
d(L, L') = |\{s : L(s) \neq L'(s)\}|
\end{equation}
with cost $C_{\text{trans}} = \lambda_m \cdot d(L_{t-1}, L_t) \cdot
\text{mc}_t$, where $\text{mc}_t$ is the per-move cost at time $t$
(including regime-specific multipliers).

\subsection{Information Regimes}

\begin{itemize}
\item \textbf{Ex-post Beam Full-Information Policy (BFIP)}: sees the
complete future trajectory $S_{1:T}$; serves as the (approximate) upper
bound. Beam search with myopic-trajectory injection guarantees
$C_{\text{BFIP}} \leq C_{\text{Myopic}}$.
\item \textbf{Receding-Horizon Warehouse Policy (RHC)}: sees only
$\text{Information}_t$ (current state, known promotions, forecasts, known
cost schedules); optimizes over a lookahead window $H$ using an internal
cost model; executes the first action. This is the deployable policy class.
\item \textbf{Myopic (ex-post)}: per-period argmin over realized costs;
clairvoyant but greedy. The trap victim.
\item \textbf{Greedy-FC}: per-period argmin over forecast-model costs; the
realistic deployable baseline.
\end{itemize}

\section{Expert Library}\label{sec:experts}

\begin{table}[h]
\centering
\caption{OR Expert Library v3.1. All experts are capacity-constrained
($\lceil n/L \rceil$ per location).}
\label{tab:experts}
\begin{tabular}{llll}
\toprule
ID & Expert & Decision rule & Best regime \\
\midrule
E1 & Static ABC & rank by full-history frequency & stable, high move cost \\
E2 & COI & rank by frequency / volume & space-constrained \\
E3 & Affinity & recency-weighted co-pick clustering & strong baskets \\
E4 & Forecast & rank by informed forecast $p_{50}$ & known promotions \\
E5 & Robust & rank by $p_{50} - \kappa \cdot (p_{90} - p_{10})$ & high uncertainty \\
E6 & DDSR-lite & opportunistic payback-gated swaps & known-order horizon \\
E7 & Joint & CP-SAT: forecast cost + move penalty & cost-coupled \\
\bottomrule
\end{tabular}
\end{table}

Experts E1--E5 produce layouts independent of the incoming layout
(cached per period). E6 (DDSR-inspired) and E7 (joint optimizer) read the
current layout: E6 swaps only when expected saving exceeds a payback
threshold; E7 solves a CP-SAT model with explicit move penalties.

\section{Experimental Environment}\label{sec:environment}

\subsection{Synthetic Platform (Primary)}

\textbf{World}: 120 SKUs (Zipf-distributed demand, category basket
structure with concentration 0.7), 60 rack locations (6m bay spacing,
$\sim$120m depth), 28-day horizon.

\textbf{Regime sequence} (8 phases): stable(4d) $\to$ promo
ramp(2d) $\to$ promo peak(4d) $\to$ promo decay(2d) $\to$ stable(3d)
$\to$ velocity reversal(6d) $\to$ affinity shift(4d) $\to$ move-cost
shock(3d, $\times$20). Promotions are \emph{known events} (multiplied into
forecasts); demand transitions ramp gradually.

\textbf{Evaluation}: per-period realized cost = L0 route distance (greedy
TSP per order) + $\lambda_m \cdot$ moves $\cdot$ mc\_unit. Beam width 30
(validated against exact enumeration, Section~\ref{sec:controlled});
deterministic CP-SAT (single worker, deterministic time budget).

\textbf{Protocol discipline}: experts slot on \emph{history} only
(known-at-time $\leq t$); evaluation on \emph{future} orders; capacity
audit on every plan; non-vacuous validation (empty fact tables rejected).

\subsection{Real Data (Supplementary)}

\textbf{Instacart} (Track B): 3.4M baskets, user-level train/eval split
(no identity leakage), top-120 SKUs by train-side frequency. Real basket
co-occurrence concentration: 0.23 (vs synthetic 0.70).

\textbf{WEPA} (external validity): real block-stacking warehouse layout
(1{,}074 storage cells, 1.4m grid) + 411{,}830 retrieval orders over 3
months, bundled with the SLAPStack package~\citep{wepastacks2024}.

\section{Expert Diversity}\label{sec:diversity}

\begin{finding}[T0: No Dominant Expert]
On the synthetic regime sequence (3 seeds $\times$ 7 evaluated phases),
the top expert (E7 Joint) wins only 52\% of periods, with 6 distinct
winners observed. Winner switches align with regime phases: promotion
$\to$ E3/E4 family, stable $\to$ E1, move-cost shock $\to$ E7 (100\%).
\end{finding}

This establishes that a fixed-expert policy is suboptimal---but does not
yet establish that \emph{sequential} routing adds value over per-period
selection. That is the subject of Section~\ref{sec:trap}.

\section{Myopic Failure: The Trap Window}\label{sec:trap}

\subsection{Definitions}

\begin{definition}[Normalized Trap Gain]
For a future window $H$ starting at period $t$:
\begin{equation}
\text{NTG}_t(H) = \frac{C^{\text{Myopic}}_{t:t+H} - C^{\text{Dynamic}}_{t:t+H}}
{C^{\text{Myopic}}_{t:t+H}}
\end{equation}
\end{definition}

\begin{definition}[Material Trap]
$\text{MaterialTrap}_t = \mathbb{1}[\text{NTG}_t > \tau]$ with $\tau = 1\%$.
\end{definition}

\subsection{Existence (T1a)}

\begin{finding}[Constructive Evidence]
On seed 17, the dynamic (BFIP) trajectory deliberately accepts a 1.6\%
current-period loss in the affinity-shift phase (choosing E1 over E3),
which positions the layout so that the subsequent move-cost-shock period
requires only 5 moves instead of 55. Total long-term gain: 5.34\%
(TrapScore = 2116/80 = 26.3).
\end{finding}

This is the core mechanism in miniature: \emph{myopic optimization
over-configures to current demand, then pays dearly to undo it when costs
spike}.

\subsection{Prevalence (T1b)}

\begin{finding}[Insurance-Shaped Distribution]
Across 12 independent seeds: 10/12 have gap $\approx$ 0, 1/12 small
(0.24--1.2\%), 1/12 large (5.34\%). Divergence points concentrate at
promo-ramp (4/12) and stable-2 (4/12)---the ``eve of structural shifts.''
Material traps: 1/12 seeds.
\end{finding}

The value of sequential routing is \emph{event-dependent}---like insurance,
it costs nothing when nothing happens and saves substantially when a trap
does occur. This motivates the controlled experiments of
Section~\ref{sec:controlled}.

\subsection{Exact Validation}\label{sec:controlled}

On small instances (60 SKU, 30 locations, 4 evaluated periods,
$7^4 = 2{,}401$ trajectories exhaustively enumerated):

\begin{finding}[Beam = Exact]
Beam-30 achieves the exact optimum on all 4 seeds (gap 0.00\%).
\textbf{With} move-cost shock ($\times$20): myopic gap 0.46--2.68\%,
trajectories diverge. \textbf{Without} shock: exact = beam = myopic
(gap 0.00\% everywhere).
\end{finding}

The no-shock control confirms the causal structure: \emph{traps are purely
the interaction of transition-period decisions with subsequent cost
shocks}---neither alone suffices.

\subsection{Controlled Trap Phase Diagram}

Systematically varying lead time $\Delta t \in \{0,1,2,4\}$ (days between
demand transition and cost shock) and shock magnitude $M_s \in
\{2,5,10,20\}$:

\begin{finding}[Intermediate Lead-Time Band]
Traps concentrate at $\Delta t = 1$: all four magnitudes are material
(NTG 1.26--1.69\%). At $\Delta t = 0$, only $M_s = 20$ is material
(1.34\%). At $\Delta t \geq 2$, no cell is material. The myopic signature:
zero moves during transition, then 2$\times$ dynamic's moves at shock
(67 vs 35 at $\Delta t=1, M_s=2$).
\end{finding}

The non-monotonicity in $\Delta t$ is itself informative: at $\Delta t=0$
the shock arrives before any reconfiguration can ``set wrong''; at
$\Delta t \geq 2$ the wrong reconfiguration has time to be amortized.
The dangerous zone is \emph{just enough lead time to act, not enough to
recover}.

\section{The Deployment Paradox}\label{sec:deployment}

\subsection{Capture Rate}

\begin{definition}[Oracle Opportunity Capture Rate]
\begin{equation}
\text{CaptureRate} = \frac{C_{\text{GreedyFC}} - C_{\text{RHC}}}
{C_{\text{GreedyFC}} - C_{\text{BFIP}}}
\end{equation}
\end{definition}

\subsection{Results}

Twelve seeds, receding-horizon $H=2$ with three internal cost models
(linear $\sum p_{50} \cdot \text{dist}$; stop-count-aware; route-aware),
schedule-aware and schedule-blind:

\begin{finding}[Capture by Conservative Models]
The crudest model (L1 linear) captures the most opportunity relative to
ex-post myopic (mean +43.4\% on seeds with headroom); the route-aware
model captures +24.6\%; the stop-aware model is unstable. \emph{Model
fidelity does not gate capture.}
\end{finding}

\begin{finding}[The Lazy Policy Paradox]
On seed 17: ex-post myopic = 38{,}106, RHC = greedyFC = 36{,}444, BFIP =
36{,}070. The deployable policies are \emph{not better forecasters}---they
are \emph{less aggressive reconfigurers}. Their imprecision acts as
implicit stickiness that avoids the trap. The true clairvoyance premium
(BFIP vs greedyFC) is only $\sim$1.0\%.
\end{finding}

This connects directly to the MTS literature's lazy-vs-aggressive
trade-off: in switching-cost environments, \emph{conservatism is
self-protecting}. Our experiments provide the warehouse-domain empirical
instance of this theoretical result.

\section{Reconfiguration Sensitivity}\label{sec:sensitivity}

Sweeping the global move-cost multiplier $\lambda_m \in \{0, 0.25, 0.5, 1,
2, 5, 10, 20, 50\}$ (3 seeds, sequence-internal shocks disabled):

\begin{finding}[Left-Low, Mid-Peak, Right-Plateau]
The dynamic-vs-myopic gap is $\sim$0.1\% at $\lambda_m \leq 0.25$ (moves
nearly free, planning unnecessary), peaks at $\lambda_m = 10$ (mean 1.26\%,
max 3.47\%), and plateaus at 1.0--1.2\% for $\lambda_m \in [20, 50]$.
\emph{At high move costs, dynamic moves more than myopic but
better}: 164 vs 152 moves at $\lambda_m = 50$.
\end{finding}

\begin{finding}[Value Per Move]
$VPM = (\Delta C_{\text{future op}} - C_{\text{move}}) / N_{\text{moves}}$.
In absolute terms (vs never-move baseline), VPM is negative and worsens
with $\lambda_m$ ($-0.8 \to -84.0$): at high costs, even optimal moves
lose vs doing nothing. In relative terms, $VPM_{\text{dynamic}} >
VPM_{\text{myopic}}$ at \emph{all} seven $\lambda_m$ levels, with the
gap widening monotonically ($+0.1 \to +5.5$).
\end{finding}

The combined interpretation: \emph{when not to reconfigure} at high
$\lambda_m$ does not mean ``never reconfigure''; it means
\textbf{reconfigure selectively}---only when expected future value exceeds
the move cost. The dynamic policy's advantage grows with stakes precisely
because it is more selective per move.

\section{Transition Cost: An Inadequate Surrogate}\label{sec:transition}

\subsection{Hidden Reconfiguration}

\begin{finding}[Switch-Count Failure]
Across 5 cost conditions $\times$ 6 seeds (36 transitions each):
False Switch rate (expert name changed, layout $\leq 2$ moves different)
= 0\%. \textbf{Hidden Reconfiguration rate} (expert name unchanged,
layout $\geq 20$ moves different) = 17--69\%. Under switch-only penalties
($\lambda_s = 20$), the rate rises to 69\% for the dynamic policy:
policies learn to ``change content without changing name'' to avoid the
indicator penalty.
\end{finding}

\subsection{Metric Properties}

\begin{proposition}[Valid Metric Space]
Both $d(L, L') = |\{s : L(s) \neq L'(s)\}|$ (Hamming) and
$d(L, L') = \sum_s \|L(s) - L'(s)\|_2$ (Euclidean sum) satisfy symmetry
and the triangle inequality on layout pairs, constituting a valid metric
space for the MTS framework.
\end{proposition}

However, four warehouse-specific phenomena violate clean MTS assumptions:
(i) asymmetric open/close pick-face costs; (ii) capacity-infeasible
transitions (hard constraint, not merely costly); (iii) batch effects
($k$ simultaneous moves share mobilization cost); (iv) swap chains
(exchanging two SKUs' locations requires a temporary third position).
We detected 663 swap pairs across expert layout pairs; each swap's true
labor cost exceeds its Hamming count by $\geq 1$ move.

These violations motivate $C_{\text{trans}} = d(L_t, L_{t+1})$ as a
\emph{generalized switching cost} rather than a pure metric.

\section{Real Warehouse Validation}\label{sec:real}

\subsection{WEPA-Natural}

40 top-retrieval SKUs, 12{,}000 real retrieval orders, real layout
geometry (1{,}074 storage cells at 1.4m grid):

\begin{finding}[No Natural Trap]
Gap = 0.00\%; the myopic trajectory is \emph{identical} to the BFIP
trajectory. Winner diversity exists (E1/E6/E7 alternate by period), but
the greedy per-period choice is already globally optimal.
\end{finding}

\subsection{WEPA-Stress}

Injecting a demand transition (628 surge orders on bottom-half SKUs,
$\sim$30\% volume increase in one period) followed by a $\times$20
move-cost shock:

\begin{finding}[Moderate Stress Still Insufficient]
Gap remains 0.00\%. The surge changes period winners but does not create
an inter-temporal trap.
\end{finding}

\subsection{External Validity Boundary}

\begin{finding}[Boundary]
Inter-temporal reconfiguration traps require demand regime changes of a
magnitude not present in steady-state warehouse operations. This bounds
trap-aware expert routing to promotion seasons, seasonal transitions, and
product-line changes---precisely the non-stationary scenarios DWERP
targets.
\end{finding}

This is consistent with practical slotting guidelines (\eg ``re-slot only
when payback is clear''): the everyday answer to ``when not to
reconfigure'' is \emph{don't bother}.

\section{Discussion}\label{sec:discussion}

\subsection{When Sequential Routing Matters}

\begin{itemize}
\item \textbf{Demand regime change}: promotion ($\times$3--10), velocity
reversal, affinity shift---sufficient magnitude to change the optimal
layout
\item \textbf{Intermediate reconfiguration cost}: $\lambda_m \sim 10$
(peak value); near-zero (free moves $\to$ plan per-period) or extreme
(nobody moves $\to$ stickiness is trivially optimal)
\item \textbf{Intermediate lead time}: $\Delta t \sim 1$ period between
demand transition and cost shock
\end{itemize}

\subsection{When It Doesn't}

\begin{itemize}
\item Steady-state operations (WEPA-Natural: gap = 0)
\item Near-zero move costs (planning unnecessary)
\item Extremely high move costs (plateau; VPM negative vs baseline)
\end{itemize}

\subsection{The Deployment Implication}

The lazy-policy paradox suggests that practical systems should
\emph{err toward stickiness} rather than sophistication. A crude forecast
model that underestimates optimization opportunities will naturally make
fewer moves---and in trap-prone environments, this is protective. This is
the warehouse-domain instance of the MTS lazy-vs-aggressive trade-off.

\subsection{Design Implications for WMS Vendors}

\begin{enumerate}
\item Measure $d(L_t, L_{t+1})$, not algorithm switches
\item Cost-aware expert selection: predict $\hat{C}(E_i | S_t)$, not
classify winners
\item Option-value hold: when uncertain, deferring a locally beneficial
reconfiguration has positive expected value
\end{enumerate}

\section{Conclusion}\label{sec:conclusion}

We formalized, demonstrated, and bounded the inter-temporal
reconfiguration trap in dynamic warehouse slotting. The core
insight---\emph{when not to reconfigure}---is the option value of deferring
locally beneficial physical changes under future uncertainty. The
deployment paradox (conservatism = protection) suggests that practical
systems need not be clairvoyant to perform well; they need to be
appropriately reluctant to act.

Three mechanisms, protected throughout this work, constitute the paper's
contribution:
\begin{enumerate}
\item Different states $\to$ different experts are optimal (T0 diversity)
\item Current optimum $\neq$ long-run optimum (T1 trap window)
\item As reconfiguration cost grows, \emph{whether to act} becomes the
decision variable (T2 sensitivity, VPM)
\end{enumerate}

\bibliographystyle{plainnat}
\begin{thebibliography}{20}

\bibitem[Borodin et al.(1992)]{borodin1992metrical}
A.~Borodin, N.~Linial, and M.~E. Saks.
\newblock An optimal on-line algorithm for metrical task system.
\newblock In \emph{STOC}, 1992.

\bibitem[Heskett(1963)]{heskett1963cube}
J.~L. Heskett.
\newblock Cube-per-order index---a key to warehouse stock location.
\newblock \emph{Transportation and Distribution Management}, 3:27--31, 1963.

\bibitem[Karimi et al.(2026)]{karimi2026keeping}
N.~Karimi, K.~Zaerpour, and R.~de~Koster.
\newblock Keeping up with changing customer demand: Adaptive data-driven
  storage and repositioning in automated warehouses.
\newblock \emph{Transportation Research Part E}, 2026.

\bibitem[Kofler et al.(2014)]{kofler2014affinity}
A.~Kofler, S.~Beham, S.~Wagner, and M.~Affenzeller.
\newblock Affinity-based slotting in warehouses with dynamic order patterns.
\newblock In \emph{Computer Aided Systems Theory}, 2014.

\bibitem[Lorenz et al.(2024a)]{lorenz2024picking}
J.~Lorenz, J.~Otto, and M.~Gendreau.
\newblock Picking operations in warehouses with dynamically arriving orders:
  How good is reoptimization?
\newblock \emph{arXiv:2409.12619}, 2024.

\bibitem[Lorenz et al.(2024b)]{lorenz2024anticipation}
J.~Lorenz, J.~Otto, and M.~Gendreau.
\newblock The value of anticipation in warehouse operations.
\newblock \emph{arXiv}, 2024.

\bibitem[Mirzaei et al.(2022)]{mirzaei2022correlated}
M.~Mirzaei, K.~Zaerpour, and R.~de~Koster.
\newblock How to benefit from order data: Correlated dispersed storage
  assignment in robotic warehouses.
\newblock \emph{International Journal of Production Research},
  60(2):549--568, 2022.

\bibitem[Waubert~de~Puiseau et al.(2022)]{waubert2022drl}
W.~Waubert~de~Puiseau, T.~Ponsen, G.~Van~Leare, and E.~Lefeber.
\newblock Dynamic storage location assignment in warehouses using deep
  reinforcement learning.
\newblock \emph{Technologies}, 10(6):129, 2022.

\bibitem[Yuan et al.(2024)]{yuan2024contextual}
Z.~Yuan et~al.
\newblock Contextual bandit policies for put-away in warehouses.
\newblock In \emph{Winter Simulation Conference}, 2024.

\bibitem[WEPAStacks(2024)]{wepastacks2024}
SLAPStack contributors.
\newblock WEPAStacks: Real warehouse layout and order data.
\newblock \url{https://pypi.org/project/slapstack/}, 2024.

\end{thebibliography}

\end{document}
